\section{Experiments}
\label{sec:experiments}

\subsection{Experimental Design, Benchmark Construction, and Reporting Protocol}
\label{sec:exp_design}

The experiments are organized around three questions that are deliberately separated in the reporting. First, in a low-dimensional shallow network, does the role-aware flow recover the same posterior structural and continuous uncertainty as a matched MCMC reference? Second, when the data-generating mechanism changes while the ambient sparsity level is held fixed, does DSS-LVR continue to identify the same small set of relevant predictors? Third, on one fixed nonlinear benchmark, how much function recovery can be retained by DSS-LVR and existing sparse-network procedures at the structural sparsity levels produced by their native selection mechanisms? Separating these questions prevents predictive fit, variable-selection accuracy, posterior fidelity, and network compression from being collapsed into a single score.

Except for the deliberately small shallow-MCMC experiment in Section~\ref{sec:shallow_exp}, all multilayer simulations use the same sample size and sparsity regime: $n=5000$, $p=200$, and $s=10$ active predictors. Thus only five percent of the available predictors are relevant. Predictors are generated independently from $\operatorname{Unif}(-2.5,2.5)$, Gaussian observation noise has variance $\sigma^2=1$, and the noise-free signal is centered and rescaled to have standard deviation $1.5$. The identities of the ten active predictors are held fixed across simulation replicates so that selected predictor sets are directly comparable, whereas the observations, teacher supports, teacher weights, nonlinear components, and noise are regenerated from the replicate seed. We use ten paired seeds, $400,\ldots,409$, and all methods fitted within one replicate receive exactly the same generated data and train--test split.

The nonlinear teacher is constructed from six sparse ridge components. Each teacher component involves exactly two active predictors, and the six two-predictor supports are generated so that all ten active predictors are represented. There are therefore twelve predictor--component memberships for ten active predictors, allowing limited repeated use of a predictor without making individual ridge functions high-dimensional. For a teacher activation $\phi$, the base signal has the form
\[
    f_0(x)=\sum_{h=1}^{6} a_h\,\phi(w_h^\top x_{\mathcal A}-b_h),
\]
where each $w_h$ has exactly two nonzero coordinates. We consider both ReLU components, $\phi(t)=(t)_+$, and trigonometric components, $\phi(t)=\sin(t)$. The latter deliberately creates model misspecification because every fitted comparison network uses ReLU hidden units.

Table~\ref{tab:selection_conditions} varies the nonlinear data-generating condition while keeping $n$, $p$, $s$, the fitted architecture, and the signal-to-noise scale fixed. Five conditions are used: ReLU; ReLU plus two pairwise interactions; ReLU plus two quadratic terms; trigonometric ridge functions; and trigonometric ridge functions plus two pairwise interactions. We do not include a combined ReLU interaction-plus-quadratic condition. Importantly, an interaction is a standardized raw predictor product $x_jx_k$, not a product of transformed trigonometric terms such as $\sin(x_j)\cos(x_k)$. In the trigonometric-interaction setting, interaction pairs are preferentially drawn from predictor pairs already appearing together in a two-predictor teacher component. This construction therefore extracts an additional low-order nonlinear interaction from the same sparse predictor structure rather than creating a second unrelated set of active variables. Quadratic additions are standardized $x_j^2$ terms. After these components are combined, the full signal is centered and rescaled to the common standard deviation of $1.5$.

All DSS-LVR multilayer fits use the same $200\rightarrow20\rightarrow20\rightarrow1$ ReLU architecture, predictor and unit primitive groups with induced edge/path states, six role-aware affine IAF layers under the cyclic three-role ordering, and the normalized-ReQU optimization map. Training uses 2000 epochs with the first 300 epochs as the all-on representation warmup. Posterior structural summaries always use the hard events $V_g>\tau_{r(g)}$ and never the continuous training gates. The architecture and optimization settings are fixed before the ten-seed evaluation and are not retuned by simulation condition. This is particularly important for the trigonometric conditions: improved or degraded performance reflects robustness to the changed response surface rather than condition-specific architecture search.

The two main tables answer different questions. Table~\ref{tab:selection_conditions} is a variable-selection table and therefore intentionally omits function MSE and $R_f^2$. It reports TPR, FPR, selection accuracy, AUROC, AUPRC, selected-support size, and repeated-fit stability. The raw MSE and $R_f^2$ values are retained as diagnostic outputs and can be reported in the appendix, but excluding them from the first table keeps the interpretation focused on predictor recovery across nonlinear conditions. Table~\ref{tab:function_sparsity}, in contrast, fixes the data-generating mechanism to the trigonometric-interaction condition and compares function recovery and structural sparsity across methods. Its main quantities are signal MSE, $R_f^2$, Kuncheva stability where predictor scores are defined, and three complementary sparsity summaries: hard retained-parameter density $D_{\mathrm{param}}$, network density $D_E$, and complete active-path density $D_\pi$.

Repeated-fit stability is evaluated throughout the experiments whenever a method produces a predictor score or ranking. Because the selected support size can vary across posterior median-probability models, we define the stability set for replicate $r$ as the top-$s$ predictors according to the method's predictor score, with $s=10$ fixed by the simulation truth. For two selected sets $S_a$ and $S_b$ of common size $s$ among $p$ candidates, the Kuncheva index is
\[
    \operatorname{KI}(S_a,S_b)=\frac{|S_a\cap S_b|p-s^2}{s(p-s)}.
\]
We report the mean over all replicate pairs. Holding the active predictor labels fixed across generated datasets makes this comparison well-defined. The resulting number measures reproducibility under the combined sampling and optimization variability of the complete fitting procedure; it should not be interpreted as an initialization-only diagnostic. For hidden-node-only methods that do not define a predictor ranking, the predictor Kuncheva index is not applicable and is reported as ``--'' rather than manufactured from another quantity.

The fixed comparison in Table~\ref{tab:function_sparsity} uses six external procedures in addition to DSS-LVR: LBBNN-LRT, LBBNN-FLOW, ISLaB-FLOW, spike-and-slab group Lasso (SS-GL), spike-and-slab group horseshoe (SS-GHS), and the input-skip $\ell_1$ network IS-ANN-L1. These methods cover connection-level latent-binary sparsity, input-skip/path sparsity, hidden-node group sparsity, and a frequentist sparse-network control. Each method is evaluated according to its native structural parameterization and sparsification rule; we do not impose a universal $0.5$ cutoff when that is not the rule used by the method itself. Where an external implementation permits the common $20\times20$ hidden architecture, we use it so that function approximation capacity is comparable. The same paired datasets are used by all methods. The weight-sharing BNN of \citet{mishra2026weightsharing} remains an important feature-selection reference but is not included in the common structural benchmark because the currently released regression workflow requires substantial adaptation and its supplied synthetic workflow uses an externally fixed top-$k$ feature step. This implementation limitation is noted below Table~\ref{tab:function_sparsity}; the method remains discussed in the Introduction.

The three sparsity summaries in Table~\ref{tab:function_sparsity} are intentionally retained together because they are not interchangeable. First,
\[
D_{\mathrm{param}}=
\frac{\#\{\text{connection weights retained in the final hard/native sparse network}\}}
{\#\{\text{candidate connection weights in the corresponding dense architecture}\}},
\]
with biases excluded from numerator and denominator. This is the common hard-network quantity available across methods. Second, $D_E$ is the method-native network or edge density. For DSS-LVR it is the posterior expected density of induced active edges,
\[
D_E=\frac{1}{|\mathcal E|}\sum_{e\in\mathcal E}P(I_e=1\mid\mathcal D),
\]
whereas deterministic or median-probability competitors contribute their corresponding native edge density. Third, $D_\pi$ is the fraction of complete active input--output computational paths; for DSS-LVR it is averaged draw-by-draw under the joint structural posterior. A model can have similar $D_{\mathrm{param}}$ or edge density to another model while retaining a very different number of complete routes from the selected predictors to the output, so $D_\pi$ captures an additional level of structural compression. None of these quantities is interpreted as hardware speed unless a physically pruned implementation is separately timed.

Training wall-clock time is recorded for the method-comparison experiment under a common hardware/software environment. Main summaries are means across the ten paired datasets; 95\% confidence intervals are retained in the result files and reported where space permits. Hyperparameters are chosen without access to test responses, and test observations are used only for final function-recovery evaluation. No method is allowed a method-specific regenerated dataset or a test-set-selected sparsity threshold.

\subsection{Structural Selection and Posterior Recovery in Shallow BNNs}
\label{sec:shallow_exp}

The shallow experiment serves a different purpose from the two multilayer tables. Its role is to determine whether DSS-LVR approximates the intended spike-and-slab posterior, not merely whether it can produce a useful sparse network. We retain the low-dimensional benchmark with $n=240$, $p=6$, two generating predictors, a two-unit ReLU teacher, a five-unit fitted network, and $\sigma^2=1$. Ten independently generated datasets are fitted by VI and by a matched multi-chain MCMC procedure using the same likelihood, prior, decoder, and structural grouping. Detailed MCMC convergence diagnostics and the numerical posterior-recovery table are moved to the appendix so that the first main-text numerical table is the larger nonlinear selection experiment.

For predictor groups, the hard posterior event for predictor $k$ controls all first-layer connections leaving that input. We report TPR, FPR, AUROC, and AUPRC from posterior inclusion probabilities, together with top-two Kuncheva stability across replicates. For unit groups, one structural state controls the incoming weights, bias, and outgoing weight of a hidden component. Since hidden units are permutation non-identifiable, recovery is summarized by the posterior active-unit count $K$ through $|E(K\mid\mathcal D)-K^\star|$ rather than by matching arbitrary fitted unit labels.

Posterior fidelity is evaluated separately for active and inactive structures. For teacher-active groups we compare conditional continuous posteriors using
\[
D_{\mathrm{SKL},A}=\operatorname{median}_{g\in\mathcal A}
 d_{\mathrm{SKL}}\{\widehat\pi_g^{\mathrm{VI}}(\theta_g\mid I_g=1),
 \widehat\pi_g^{\mathrm{MCMC}}(\theta_g\mid I_g=1)\},
\]
and for teacher-inactive groups we compare inclusion uncertainty using
\[
D_{\mathrm{JS},0}=\operatorname{median}_{g\notin\mathcal A}
 d_{\mathrm{JS}}\{\operatorname{Bern}(\widehat{\mathrm{PIP}}_g^{\mathrm{VI}}),
 \operatorname{Bern}(\widehat{\mathrm{PIP}}_g^{\mathrm{MCMC}})\}.
\]
A joint posterior-density figure for representative generating components is retained in the main text because it directly shows whether the role-aware IAF captures posterior dependence that is invisible in marginal support-recovery summaries. Function MSE and $R_f^2$ remain supplementary shallow diagnostics, while the numerical posterior-recovery table is reported in the appendix.

\subsection{Predictor Recovery Across Nonlinear Simulation Conditions}
\label{sec:selection_conditions}

We next hold the dimensionality and sparsity regime fixed at $n=5000$, $p=200$, and $s=10$ and vary only the nonlinear construction described in Section~\ref{sec:exp_design}. This experiment tests whether predictor semantics survive changes in the generating response surface when the fitted DSS-LVR architecture and training schedule remain unchanged. The ReLU conditions provide settings relatively close to the fitted activation family, whereas the trigonometric conditions require a ReLU MLP to approximate smooth oscillatory ridge functions. The interaction settings add two standardized raw products of active predictors. Because the trigonometric teacher already assigns two predictors to each sparse component, these product terms can be drawn from the same local predictor pairs, yielding a mixed nonlinear signal without expanding the true predictor support.

The main table contains only selection quantities. Predictor $j$ is selected in the median-probability model when $\widehat{\operatorname{PIP}}_j>0.5$. We report TPR and FPR at this hard decision, overall selection accuracy, AUROC and AUPRC from the complete PIP ranking, the mean selected-support size, and the pairwise top-ten Kuncheva index. AUROC and AUPRC measure ranking quality even when a posterior is slightly over- or under-confident around the $0.5$ decision threshold; AUPRC is especially informative because only 10 of 200 predictors are active. The selected-support size prevents a method from obtaining high TPR by simply retaining many predictors, while Kuncheva distinguishes reproducible rankings from averages that hide replicate-to-replicate instability.

\begin{table}[t]
\centering
\small
\caption{Predictor recovery across nonlinear simulation conditions. All settings use $n=5000$, $p=200$, and $s=10$ with the same fitted $20\times20$ DSS-LVR MLP. MSE and $R_f^2$ are intentionally omitted from this selection-focused table.}
\label{tab:selection_conditions}
\begin{tabular}{lccccccc}
\toprule
Condition & TPR & FPR & Accuracy & AUROC & AUPRC & $|\widehat S|$ & Kuncheva \\
\midrule
ReLU & & & & & & & \\
ReLU + interaction & & & & & & & \\
ReLU + quadratic & & & & & & & \\
Trigonometric & & & & & & & \\
Trigonometric + interaction & & & & & & & \\
\bottomrule
\end{tabular}
\end{table}

A companion figure displays replicate-level predictor selection frequencies for the five conditions. Function-recovery metrics are still computed for quality control and retained in the supplementary results, but they are intentionally separated from Table~\ref{tab:selection_conditions}; the dedicated method comparison below is where function recovery is evaluated as a primary endpoint.

\subsection{Function Recovery and Sparse-Network Structure}
\label{sec:function_sparsity}

The second main experiment fixes the most representative misspecified setting from the previous experiment: the trigonometric ridge teacher with six two-predictor components and two raw pairwise interactions. The data size, true support, noise variance, and signal scale remain $n=5000$, $p=200$, $s=10$, $\sigma^2=1$, and $\operatorname{sd}\{f^\star(X)\}=1.5$. This condition is nonlinear in two distinct ways---oscillatory sparse ridge functions and low-order predictor products---while preserving a known ten-variable truth. It therefore provides a common test of whether a procedure can retain a complex function after producing a sparse neural representation.

DSS-LVR is compared with LBBNN-LRT, LBBNN-FLOW, ISLaB-FLOW, SS-GL, SS-GHS, and IS-ANN-L1. LBBNN-LRT and LBBNN-FLOW place latent binary states on neural connections, ISLaB-FLOW augments this construction with input skips and explicit active-path structure, SS-GL and SS-GHS place spike-and-slab group shrinkage on hidden nodes, and IS-ANN-L1 supplies a non-Bayesian input-skip sparse-network control. These methods do not define sparsity at the same primitive level, which is precisely why the final hard parameter fraction, native edge density, and complete path density are all retained in the comparison. Predictor Kuncheva is reported for methods with a native input score/ranking and is not assigned to the node-only SS-GL/SS-GHS procedures.

Function recovery is measured against the known noise-free test signal,
\[
\mathrm{MSE}_f=\frac{1}{n_{\mathrm{test}}}\sum_i\{\widehat f(X_i)-f^\star(X_i)\}^2,
\qquad
R_f^2=1-\frac{\sum_i\{\widehat f(X_i)-f^\star(X_i)\}^2}
{\sum_i\{f^\star(X_i)-\bar f^\star\}^2}.
\]
These metrics are paired with $D_{\mathrm{param}}$, $D_E$, and $D_\pi$ rather than with one generic ``sparsity'' column. For DSS-LVR, the hard MPM parameter fraction and posterior expected induced-edge density need not be identical because one summarizes a single hard posterior-summary network whereas the other averages induced edge activity over posterior draws. Path density can be substantially smaller again because complete paths require simultaneous activation of a predictor and every unit along the route. Reporting all three quantities makes this distinction explicit rather than treating them as interchangeable versions of the same metric.

\begin{table*}[t]
\centering
\scriptsize
\setlength{\tabcolsep}{4pt}
\caption{Function recovery and structural sparsity on the fixed trigonometric-interaction benchmark ($n=5000$, $p=200$, $s=10$). Lower MSE$_f$ and higher $R_f^2$ indicate better function recovery; lower $D_{\mathrm{param}}$, $D_E$, and $D_\pi$ indicate greater structural compression and must be interpreted jointly with function recovery. Higher Kuncheva indicates greater selection stability.}
\label{tab:function_sparsity}
\begin{tabular}{lccccccc}
\toprule
Method & MSE$_f$ & $R_f^2$ & Kuncheva & $D_{\mathrm{param}}$ & $D_E$ & $D_\pi$ & Time (s) \\
\midrule
DSS-LVR & & & & & & & \\
LBBNN-LRT & & & & & & & \\
LBBNN-FLOW & & & & & & & \\
ISLaB-FLOW & & & & & & & \\
SS-GL & & & -- & & & & \\
SS-GHS & & & -- & & & & \\
IS-ANN-L1 & & & & & & & \\
\bottomrule
\end{tabular}
\end{table*}

The weight-sharing BNN of \citet{mishra2026weightsharing} is not included as a row in Table~\ref{tab:function_sparsity}. Public regression code is available, but the released synthetic regression workflow is not a self-contained non-oracle implementation of this common structural benchmark without substantive adaptation: it contains author-specific paths and its supplied simulation workflow ranks first-layer inclusion quantities and then explicitly selects a fixed top-$k$ set before retraining. We therefore retain wsBNN as a closely related feature-selection method in the Introduction and discuss this reproducibility limitation here rather than silently replacing its native workflow with our own implementation.

A companion figure plots MSE$_f$ against $D_{\mathrm{param}}$ for the seven methods, with $D_E$ and $D_\pi$ reported in the table. This visualization emphasizes that a sparse network is useful only if structural removal preserves the represented function; the numerical table retains the richer distinction between hard parameter compression, edge-level sparsity, and complete path sparsity.

\subsection{Ablations and Computational Controls}
\label{sec:ablation}

The ablations isolate components of the final DSS-LVR posterior rather than historical architectures that are no longer part of the proposed method. A mean-field DSS-LVR posterior controls for the value of dependent posterior transport. A generic affine IAF with coordinate-level orderings is compared with the role-aware ordering to isolate the contribution of the $U$--$V$--$\tau$ structure. Smooth exact-zero and normalized-ReQU maps are compared while keeping the hard posterior event unchanged, and IAF depth is varied to verify that conclusions do not depend on one arbitrary transport depth. Attention-conditioned coupling architectures are not included because they are alternative flow architectures rather than removable components of the final role-aware IAF.

Training wall-clock time and posterior-sampling time are recorded under a common environment, and flow parameter count is reported in the ablation. Raw per-replicate outputs are retained only for quantities needed to reproduce the tables and figures; epoch histories and large diagnostic objects are not written by the main experiment scripts. Additional MCMC diagnostics, posterior predictive coverage, condition-specific MSE/$R_f^2$ from the first multilayer experiment, and full confidence-interval summaries are reported in the appendix.
